\pagestyle{appendix}
\appendix
\section*{Appendices}
\addcontentsline{toc}{section}{Appendices}
\renewcommand{\thesubsection}{\Alph{subsection}}
%%%%%%%%%%%%%%%%%%%%%%%%%%%%%%%%%%%%%%%%%%%%%%%%

% Delete before submitting the report
\subsection{Style Guidelines}

\subsubsection{Images and Figures}

\begin{figure}[htb]
  \centering
  \includegraphics[width=\linewidth]{example-image-a}
  \caption{This is a sample image.}
  \label{fig:sample}
\end{figure}

\subsubsection{Tables}
\begin{table}[h]
  \centering
  \caption{This is a sample table.}
  \label{tab:sample}
  \begin{tabular}{lcr}
    \hline
    Column1 & Column2 & Column3 \\
    \hline
    1 & 2 & 3 \\
    \hline
  \end{tabular}
\end{table}

\subsubsection{Equations}
An example of an equation can be found in \eqref{eq:sample}. Treat equations like phrases and use punctuation. Present the equation first, followed by the definitions of its variables.

\begin{equation}
{x}_{k + 1} = {A} {x}_k + {B} {u}_k, 
\label{eq:sample}
\end{equation}

where ${x}_{k}$ denotes the system state at time $k$, ${u}_k$ represents the control variable, $A$ is the system matrix, and $B$ is the input matrix.

\subsubsection{Algorithms}

\begin{algorithm}
\caption{An algorithm with caption}\label{alg:cap}
\begin{algorithmic}
\Require $n \geq 0$
\Ensure $y = x^n$
\State $y \gets 1$
\State $X \gets x$
\State $N \gets n$
\While{$N \neq 0$}
\If{$N$ is even}
    \State $X \gets X \times X$
    \State $N \gets \frac{N}{2}$  \Comment{This is a comment}
\ElsIf{$N$ is odd}
    \State $y \gets y \times X$
    \State $N \gets N - 1$
\EndIf
\EndWhile
\end{algorithmic}
\end{algorithm}

\subsubsection{References}
You can use the \texttt{refs.bib} BibTex file to store your references. Use the command \texttt{\textbackslash cite\{reference-label\}} to cite, e.g., \cite{mlp-universal} or \cite{mlp-universal, rbf-universal}.

\newpage
\subsection{Second Appendix}

